\chapter{Singleton-Unabhängigkeit: direkte Gegenbeispiel-Variante}
\label{ch:singl-ci-direkt}

\section{Idee}
\label{sec:singl-direkt-idee}

Dieses Kapitel behandelt dasselbe Entscheidungsproblem
\(\textit{SingInd}_{\mathrm{st}}\) wie Kapitel~\ref{ch:singl-ci}, wählt aber
einen unmittelbar an der Definition der Unabhängigkeit orientierten Zugang.
Statt die Lösungsstruktur abstrakt über \emph{Produktrechtecke}
(Kapitel~\ref{ch:singl-ci}) zu beschreiben, argumentieren wir direkt mit den
drei Labellings \(\lambda_1,\lambda_2,\lambda_3\) aus
Definition~\ref{def:BU:AFs} und der Frage, ob sich zu einem Paar
\(\lambda_1,\lambda_2\) stets ein ``kreuzendes'' drittes Labelling
\(\lambda_3\) finden lässt.

Beide Kapitel beschreiben am Ende denselben Algorithmus und führen zu
denselben SAT-Aufrufen; sie unterscheiden sich nur in der Art, wie die
Korrektheit begründet wird. Abschnitt~\ref{sec:singl-direkt-vergleich} stellt
die beiden Sichtweisen explizit gegenüber. Wir verwenden durchgehend
\(L_{\mathrm{st}}=\{\mathit{in},\mathit{out}\}\), \(X=\{x\}\), \(Y=\{y\}\) mit
\(x\neq y\) und \(Z\subseteq A\setminus\{x,y\}\).

\section{Gegenbeispiele im Singleton-Fall}
\label{sec:singl-direkt-gegenbsp}

Wir erinnern an die Definition der bedingten Unabhängigkeit
(Definition~\ref{def:BU:AFs}): \(\{x\}\bot_{\mathrm{st}}\{y\}\mid Z\) bedeutet,
dass zu \emph{je} zwei \textit{stable}-Labellings \(\lambda_1,\lambda_2\), die
auf \(Z\) übereinstimmen, ein drittes \textit{stable}-Labelling \(\lambda_3\)
existiert, das auf \(x\) mit \(\lambda_1\), auf \(y\) mit \(\lambda_2\) und auf
\(Z\) mit beiden übereinstimmt. Die Verneinung dieser Aussage führt
unmittelbar auf den Begriff des Gegenbeispiels.

\begin{Definition}
    \label{def:singl-direkt-gegenbsp}
    Sei \(F=(A,R)\) ein AF, seien \(x,y\in A\) verschieden und sei
    \(Z\subseteq A\setminus\{x,y\}\). Ein \emph{Gegenbeispiel} zu
    \(\{x\}\bot_{\mathrm{st}}\{y\}\mid Z\) ist ein Paar
    \(\lambda_1,\lambda_2\in\Lambda_{\mathrm{st}}(F)\) mit
    \begin{enumerate}
        \item \(\lambda_1|_Z=\lambda_2|_Z\) und
        \item es existiert kein \(\lambda_3\in\Lambda_{\mathrm{st}}(F)\) mit
              \[
                  \lambda_3(x)=\lambda_1(x),\qquad
                  \lambda_3(y)=\lambda_2(y),\qquad
                  \lambda_3|_Z=\lambda_1|_Z.
              \]
    \end{enumerate}
    Das in Bedingung~2 gesuchte Labelling \(\lambda_3\) nennen wir die zu
    \((\lambda_1,\lambda_2)\) gehörige \emph{Kreuzkombination}.
\end{Definition}

Anschaulich übernimmt die Kreuzkombination das \(x\)-Label vom ersten und das
\(y\)-Label vom zweiten Labelling, während die Bedingung \(Z\) unverändert
bleibt. Ein Gegenbeispiel ist also genau ein auf \(Z\) übereinstimmendes Paar,
dessen Kreuzkombination sich nicht zu einem \textit{stable}-Labelling
ergänzen lässt.

\begin{Proposition}
    \label{prop:singl-direkt}
    Es gilt \(\{x\}\bot_{\mathrm{st}}\{y\}\mid Z\) genau dann, wenn kein
    Gegenbeispiel im Sinne von Definition~\ref{def:singl-direkt-gegenbsp}
    existiert.
\end{Proposition}

\begin{proof}
    \label{pr:singl-direkt}
    Dies ist die wörtliche Verneinung von Definition~\ref{def:BU:AFs} im Fall
    \(X=\{x\}\), \(Y=\{y\}\). Unabhängigkeit besagt, dass zu jedem auf \(Z\)
    übereinstimmenden Paar \(\lambda_1,\lambda_2\) eine Kreuzkombination
    existiert. Ein Gegenbeispiel ist genau ein Paar, für das dies fehlschlägt.
    Folglich gilt Unabhängigkeit genau dann, wenn es kein solches Paar gibt.
\end{proof}

Die folgende Beobachtung schränkt die Suche stark ein: In einem Gegenbeispiel
müssen sich \(\lambda_1\) und \(\lambda_2\) sowohl auf \(x\) als auch auf \(y\)
unterscheiden. Stimmen sie auf einem der beiden Argumente überein, so ist die
Kreuzkombination trivial vorhanden.

\newpage

\begin{Lemma}
    \label{lem:singl-direkt-opposite}
    Ist \((\lambda_1,\lambda_2)\) ein Gegenbeispiel, so gilt
    \(\lambda_1(x)\neq\lambda_2(x)\) und \(\lambda_1(y)\neq\lambda_2(y)\).
\end{Lemma}

\begin{proof}
    \label{pr:singl-direkt-opposite}
    Angenommen \(\lambda_1(x)=\lambda_2(x)\). Dann erfüllt
    \(\lambda_3:=\lambda_2\) bereits alle Bedingungen an die
    Kreuzkombination, denn
    \[
        \lambda_3(x)=\lambda_2(x)=\lambda_1(x),\quad
        \lambda_3(y)=\lambda_2(y),\quad
        \lambda_3|_Z=\lambda_2|_Z=\lambda_1|_Z.
    \]
    Damit wäre \((\lambda_1,\lambda_2)\) kein Gegenbeispiel, im Widerspruch
    zur Annahme. Also \(\lambda_1(x)\neq\lambda_2(x)\). Der Fall für \(y\)
    folgt symmetrisch mit \(\lambda_3:=\lambda_1\).
\end{proof}

Da \(L_{\mathrm{st}}\) nur die beiden Werte \(\mathit{in}\) und
\(\mathit{out}\) besitzt, legt Lemma~\ref{lem:singl-direkt-opposite} die Labels
von \(x\) und \(y\) in einem Gegenbeispiel bis auf Vertauschung von
\(\lambda_1\) und \(\lambda_2\) auf genau zwei \emph{Konfigurationen} fest:
\[
\begin{aligned}
    \text{(K1)}\quad &
    \lambda_1(x)=\mathit{in},\ \lambda_1(y)=\mathit{in},
    &\quad
    \lambda_2(x)=\mathit{out},\ \lambda_2(y)=\mathit{out};\\
    \text{(K2)}\quad &
    \lambda_1(x)=\mathit{in},\ \lambda_1(y)=\mathit{out},
    &\quad
    \lambda_2(x)=\mathit{out},\ \lambda_2(y)=\mathit{in}.
\end{aligned}
\]
In Konfiguration~(K1) lautet die gesuchte Kreuzkombination
\(\lambda_3(x)=\mathit{in}\), \(\lambda_3(y)=\mathit{out}\); in~(K2) ist es
\(\lambda_3(x)=\mathit{in}\), \(\lambda_3(y)=\mathit{in}\). Da die Definition
der Unabhängigkeit über \emph{alle} geordneten Paare \((\lambda_1,\lambda_2)\)
quantifiziert, muss zusätzlich die durch Vertauschen von \(\lambda_1\) und
\(\lambda_2\) entstehende Kreuzkombination realisierbar sein --- in~(K1) also
auch \(\lambda_3(x)=\mathit{out}\), \(\lambda_3(y)=\mathit{in}\). Ein
Gegenbeispiel liegt vor, sobald für eine gemeinsame \(Z\)-Belegung eine dieser
beiden Kreuzkombinationen fehlt.

\section{Algorithmus}
\label{sec:singl-direkt-algo}

\subsection{Realisierbarkeitsabfragen}
\label{subsec:singl-direkt-realize}

Das Verfahren stützt sich auf eine einzige
Hilfsabfrage. Für ein Argument \(a\in A\) und ein Label
\(\ell\in L_{\mathrm{st}}\) schreiben wir \(\mathrm{Fix}_a(\ell)\) für das
Literal, das \(a\) auf \(\ell\) festlegt, also \(\mathrm{Fix}_a(\mathit{in}):=
i_a\) und \(\mathrm{Fix}_a(\mathit{out}):=\neg i_a\) mit der \textit{in}-Variable
\(i_a\) aus Definition~\ref{def:st-pl}. Für eine Teilbelegung
\(\pi:W\to L_{\mathrm{st}}\) mit \(W\subseteq A\) sei
\[
    \mathrm{realisierbar}(\pi)
    \quad\Longleftrightarrow\quad
    \mathrm{st}_F(\mathcal L)\land
    \bigwedge_{a\in W}\mathrm{Fix}_a(\pi(a))
    \text{ ist erfüllbar.}
\]
Nach Lemma~\ref{lem:st-pl-correct} gilt \(\mathrm{realisierbar}(\pi)\) genau
dann, wenn ein \(\lambda\in\Lambda_{\mathrm{st}}(F)\) mit \(\lambda|_W=\pi\)
existiert. Für die Aufzählung gemeinsamer \(Z\)-Belegungen verwenden wir
zwei Kopien \(\mathcal L^1,\mathcal L^2\) der Labelling-Variablen mit den
zugehörigen Literalen \(\mathrm{Fix}^1_a(\ell)\), \(\mathrm{Fix}^2_a(\ell)\)
und der Gleichheitsformel \(\mathrm{Eq}^{1,2}_Z\) aus
Definition~\ref{def:eq-qbf}.

\subsection{Verfahren}
\label{subsec:singl-direkt-verfahren}

Der Algorithmus sucht ein Gegenbeispiel, indem er für jede der beiden
Konfigurationen die gemeinsamen \(Z\)-Belegungen aufzählt und jeweils prüft,
ob eine Kreuzkombination fehlt.

\begin{enumerate}
    \item \textbf{Globale Erfüllbarkeit.} Ist
    \(\mathrm{realisierbar}(\emptyset)\) falsch, so gilt
    \(\Lambda_{\mathrm{st}}(F)=\emptyset\). Dann gibt es keine zwei Labellings
    \(\lambda_1,\lambda_2\), also kein Gegenbeispiel, und
    \(\{x\}\bot_{\mathrm{st}}\{y\}\mid Z\) ist vakuos erfüllt.

    \item \textbf{Fixed-Label-Vortest.} Prüfe für \(a\in\{x,y\}\) und
    \(\ell\in L_{\mathrm{st}}\) jeweils \(\mathrm{realisierbar}(a\mapsto\ell)\).
    Nimmt \(x\) oder \(y\) unter allen \textit{stable}-Labellings nur ein
    einziges Label an, so ist \(\lambda_1(x)\neq\lambda_2(x)\) (bzw.
    \(\lambda_1(y)\neq\lambda_2(y)\)) unmöglich. Nach
    Lemma~\ref{lem:singl-direkt-opposite} kann dann kein Gegenbeispiel
    existieren, und es gilt \(\{x\}\bot_{\mathrm{st}}\{y\}\mid Z\).

    \item \textbf{Enumeration der Konfigurationen.} Für jede Konfiguration
    \[
        K=\bigl((\alpha_1,\beta_1),(\alpha_2,\beta_2)\bigr)\in\{\text{(K1)},\text{(K2)}\}
    \]
    mit den oben festgelegten Labels:
    \begin{enumerate}
        \item Setze \(B:=\top\) und bilde die Formel, die
        \(\lambda_1,\lambda_2\) in die Konfiguration \(K\) zwingt und ihre
        Gleichheit auf \(Z\) erzwingt:
        \[
        \begin{aligned}
            C^{12}_{K}
            :=\ &
            \mathrm{st}_F(\mathcal L^1)\land\mathrm{st}_F(\mathcal L^2)
            \land\mathrm{Fix}^1_x(\alpha_1)\land\mathrm{Fix}^1_y(\beta_1)\\
            &\land\mathrm{Fix}^2_x(\alpha_2)\land\mathrm{Fix}^2_y(\beta_2)
            \land\mathrm{Eq}^{1,2}_Z.
        \end{aligned}
        \]
        \item \textbf{Gemeinsame \(Z\)-Belegung suchen.} Solange
        \(C^{12}_{K}\land B\) erfüllbar ist, wähle ein Modell und lies die
        gemeinsame Belegung \(\eta=\lambda_1|_Z=\lambda_2|_Z\) ab. 
        \item \textbf{Kreuzkombinationen prüfen.} Ist
        \[
            \neg\,\mathrm{realisierbar}\bigl(\{x\mapsto\alpha_1,\,y\mapsto\beta_2\}\cup\eta\bigr)
            \ \text{ oder }\
            \neg\,\mathrm{realisierbar}\bigl(\{x\mapsto\alpha_2,\,y\mapsto\beta_1\}\cup\eta\bigr),
        \]
        so fehlt eine der beiden Kreuzkombinationen. Damit ist
        \((\lambda_1,\lambda_2)\) ein Gegenbeispiel; der Algorithmus gibt es
        zurück und terminiert mit \(\{x\}\not\bot_{\mathrm{st}}\{y\}\mid Z\).
        \item \textbf{Blocking.} Sind beide Kreuzkombinationen realisierbar,
        so ist \(\eta\) für diese Konfiguration kein Gegenbeispiel. Schließe
        \(\eta\) durch die Blocking-Klausel
        \[
            \mathrm{Block}(\eta):=\bigvee_{z\in Z}\neg\,\mathrm{Fix}^1_z(\eta(z))
        \]
        aus, setze \(B:=B\land\mathrm{Block}(\eta)\) und fahre mit der
        nächsten gemeinsamen \(Z\)-Belegung fort.
    \end{enumerate}

    \item \textbf{Abschluss.} Wurden beide Konfigurationen ohne Gegenbeispiel
    durchlaufen, so existiert nach Lemma~\ref{lem:singl-direkt-opposite}
    überhaupt kein Gegenbeispiel.
\end{enumerate}

Wir verwenden \(\bigvee\emptyset=\bot\); im Fall
\(Z=\emptyset\) bricht daher jede Konfiguration nach genau einer Iteration ab.

\subsection{Korrektheit und Terminierung}
\label{subsec:singl-direkt-korrekt}

\begin{Lemma}
    \label{lem:singl-direkt-korrekt}
    Der Algorithmus aus Abschnitt~\ref{subsec:singl-direkt-verfahren}
    entscheidet \(\textit{SingInd}_{\mathrm{st}}\) korrekt und terminiert.
\end{Lemma}

\begin{proof}
    \label{pr:singl-direkt-korrekt}
    \emph{Korrektheit bei Gegenbeispiel.} Gibt der Algorithmus in Schritt~3(c)
    ein Paar zurück, so existieren nach Schritt~3(b) zwei
    \textit{stable}-Labellings \(\lambda_1,\lambda_2\) mit \(\lambda_1|_Z=
    \lambda_2|_Z=\eta\) in Konfiguration \(K\), und eine der beiden
    Kreuzkombinationen ist unrealisierbar. Damit existiert kein passendes
    \(\lambda_3\), das Paar ist ein Gegenbeispiel nach
    Definition~\ref{def:singl-direkt-gegenbsp}, und mit
    Proposition~\ref{prop:singl-direkt} folgt
    \(\{x\}\not\bot_{\mathrm{st}}\{y\}\mid Z\).

    \emph{Korrektheit bei Unabhängigkeit.} Meldet der Algorithmus
    Unabhängigkeit, so liegt einer von drei Fällen vor:
    \(\Lambda_{\mathrm{st}}(F)=\emptyset\) (Schritt~1), ein fixiertes Label
    (Schritt~2, Lemma~\ref{lem:singl-direkt-opposite}) oder der erschöpfende
    Durchlauf beider Konfigurationen (Schritt~4). In allen Fällen existiert
    kein Gegenbeispiel, sodass Proposition~\ref{prop:singl-direkt}
    \(\{x\}\bot_{\mathrm{st}}\{y\}\mid Z\) liefert.

    \emph{Vollständigkeit.} Gilt \(\{x\}\not\bot_{\mathrm{st}}\{y\}\mid Z\), so
    existiert nach Proposition~\ref{prop:singl-direkt} ein Gegenbeispiel
    \((\lambda_1,\lambda_2)\). Nach Lemma~\ref{lem:singl-direkt-opposite}
    unterscheiden sich \(\lambda_1,\lambda_2\) auf \(x\) und auf \(y\); sie
    realisieren also (bis auf Vertauschung) eine der Konfigurationen \(K\) mit
    gemeinsamer Belegung \(\eta^\ast:=\lambda_1|_Z\). Im Durchlauf dieser
    Konfiguration erfüllt \(\eta^\ast\) die Formel \(C^{12}_{K}\); geblockt
    werden ausschließlich Belegungen, deren Kreuzkombinationen beide
    realisierbar sind, also nicht \(\eta^\ast\). Die Enumeration erreicht somit
    \(\eta^\ast\) (oder zuvor ein anderes Gegenbeispiel), und die Prüfung in
    Schritt~3(c) schlägt fehl.

    \emph{Terminierung.} Jede Iteration der Enumerationsschleife terminiert
    entweder oder schließt durch \(\mathrm{Block}(\eta)\) genau eine
    realisierbare gemeinsame \(Z\)-Belegung dauerhaft aus. Davon gibt es
    höchstens \(2^{|Z|}\); die äußere Schleife über die beiden Konfigurationen
    ist endlich.
\end{proof}

Die Heuristiken zur Beschleunigung der Gegenbeispielsuche --- inkrementelles
SAT, Blocking über UNSAT-Cores, verallgemeinertes Blocking, das globale
Backbone von \(Z\), SCC-Vorverarbeitung sowie eine günstige Reihenfolge der
Abfragen --- gelten unverändert; sie sind in
Abschnitt~\ref{subsec:singl-heuristiken} beschrieben.

\section{Vergleich beider Charakterisierungen}
\label{sec:singl-direkt-vergleich}

Beide Kapitel entscheiden \(\textit{SingInd}_{\mathrm{st}}\) mit demselben
Verfahren; sie unterscheiden sich allein in der theoretischen Begründung.

\begin{itemize}
    \item \textbf{Produktrechteck-Sicht (Kapitel~\ref{ch:singl-ci}).} Für jede
    realisierbare \(Z\)-Belegung \(\eta\) wird die Menge \(R_\eta\) der unter
    \(\eta\) möglichen \((x,y)\)-Labelpaare betrachtet. Unabhängigkeit gilt
    genau dann, wenn jedes \(R_\eta\) ein Produktrechteck \(P\times Q\) ist.
    Diese Sicht ist kompakt und macht die kombinatorische Struktur explizit,
    ist aber von der ursprünglichen Definition abstrahiert.

    \item \textbf{Gegenbeispiel-Sicht (dieses Kapitel).} Es wird direkt nach
    einem Paar \((\lambda_1,\lambda_2)\) ohne Kreuzkombination gesucht. Diese
    Sicht liest sich unmittelbar aus Definition~\ref{def:BU:AFs} ab und ist
    näher an der Implementierung, benennt die kombinatorische Gesamtstruktur
    aber nicht explizit.
\end{itemize}

Die beiden Sichten sind äquivalent und übersetzen sich Schritt für Schritt
ineinander. Eine \emph{Diagonale} \(D_1=\{(\mathit{in},\mathit{in}),
(\mathit{out},\mathit{out})\}\) bzw. \(D_2=\{(\mathit{in},\mathit{out}),
(\mathit{out},\mathit{in})\}\), der in \(R_\eta\) ein Kreuzpaar fehlt,
entspricht genau einer Konfiguration~(K1) bzw.~(K2), deren Kreuzkombination
unrealisierbar ist. Eine fehlende Rechteck-Ecke \((\alpha_1,\beta_2)\notin
R_\eta\) ist dasselbe wie eine fehlende Kreuzkombination \(\lambda_3\) mit
\(\lambda_3(x)=\alpha_1\), \(\lambda_3(y)=\beta_2\). Folglich erzeugen beide
Algorithmen für jede Eingabe dieselbe Folge von Realisierbarkeitsabfragen und
haben dieselbe Aufrufkomplexität (vgl.
Korollar~\ref{cor:singl-aufrufe}).
